\section{Related work}


In 2009, Krishnamurthy and Wills provided one of the first longitudinal analyses of user information flows to third-party sites (called \textit{aggregators})~\cite{krishnamurthy2009privacy}.
The authors also observed a trend of serving third-party tracking content from first-party contexts, pointing out the challenges for countermeasures based on blocklists.
Meyer and Mitchell studied the technology and policy aspects of third-party tracking~\cite{mayer2012third}. Englehardt and Narayanan \cite{englehardt2016online} 
measured tracking on Alexa top million websites using OpenWPM and discovered new fingerprinting techniques such as AudioContext API-based fingerprinting.

The CNAME tracking scheme was mentioned anecdotally by Bau in 2013 \cite{bau2013promising}, but the authors did not focus on the technique specifically.
To our knowledge, the first systematic analysis of the CNAME scheme used to embed third-party trackers in first-party content is the work of Olejnik and Casteluccia \cite{olejnik2014analysis}, 
in which they identified this special arrangement as part of the real-time bidding setup.
The authors also reported leaks of first-party cookies to such third parties. In our paper, we extensively expand such analyses.
Although cookies were most commonly used for cross-site tracking, more advanced mechanisms have been deployed by websites and studied by the researchers. Browser fingerprinting \cite{eckersley2010unique}, where traits of the host \cite{yen2012host}, system, browser and graphics stack~\cite{mowery2012pixel} are extracted to identify the user is one of the stateless tracking vectors that does not need cookies to operate.
Fingerprinting on the web was measured at scale by Acar et al. \cite{acar2013fpdetective, Acar:2014:WNF:2660267.2660347}, Nikiforakis et al.\cite{nikiforakis2013cookieless},
and Englehardt and Narayanan~\cite{englehardt2016online}. Combining multiple tracking vectors at the same time may give rise to supercookies or evercookies, as demonstrated first by Samy Kamkar \cite{kamkar2010evercookie}.
Over the years, many information exfiltration or tracking vectors have been studied, including Cache Etag HTTP header~\cite{ayenson2011flash}, WebSockets \cite{bashir2018tracking}, ultrasound beacons \cite{mavroudis2017privacy}, and fingerprinting sensors calibrations on mobile devices with sensors \cite{zhang2019sensorid}. 
 

Similar to these studies we measure the prevalence of a tracking mechanism that tries to circumvent existing countermeasures. However our work uses novel methods to identify CNAME-based trackers in historical crawl data, allowing us to perform a longitudinal measurement.


In concurrent work, Dao et al.\ also explored the ecosystem of CNAME-based trackers~\cite{dao2020cnamecloaking}.
Based on a crawl of the Alexa top 300k, they find 1,762 CNAME-based tracking domains as of January 2020, which are detected by matching the CNAME domain with EasyPrivacy.
In our work, we detected 9,273 sites that leverage CNAME-based tracking in a same-site context and an additional 19,226 websites that use it in a cross-site context.
We rely on an approach that combines historical DNS records (A records) with manually constructed fingerprints.
The latter is used to filter out any potential false positives that may be caused by changes in the IP space ownership, or because the CNAME- or A-records may be used to other services of the same provider unrelated to tracking.
Based on the evaluation of our method in Section~\ref{sec:method_validation}, we find that it is important to use request-specific information to prevent incorrectly marking domains as using CNAME-based tracking.
Furthermore, relying on filter lists, and in particular on the eTLD+1 domains that are listed, could result in the inclusion of non-tracking domains, e.g.\ sp-prod.net is the second most popular tracker considered by Dao et al., but was excluded in our work as it is part of a ``Consent Management Platform'' that captures cookie consent for compliance with GDPR \cite{CMP}.
Additionally, filter lists may be incomplete, resulting in trackers being missed: for example, Pardot, the tracker we find to be most widely used, was not detected in prior work.
Consequently, relying on filter lists also prevents the detection of new trackers, this limitation is not applicable to our method.

Dao et al.\ also perform an analysis of the historical evolution of CNAME-based tracking, based on four datasets of the Alexa top 100k websites collected between January 2016 and January 2020.
As the used OpenWPM datasets do not include DNS records, the researchers rely on a historical forward DNS dataset provided by Rapid7~\cite{rapid7fdns}, which does not cover all domains over time.
By using the HTTP Archive dataset, which includes the IP address that was used, we were able to perform a more granular analysis, showing a more accurate growth pattern.
We also show that this growth is rapidly increasing, significantly outperforming third-party trackers with a comparable customer base.
Finally, to the best of our knowledge, we are the first to perform an analysis of the privacy and security implications associated with the CNAME-based tracking scheme.
